\scanpage{41}
\margin{40}
avoir, par le procédé incomplet précédent, en \unsure{recollant} rien le long des arcs de spécialisation exceptionnels, une \uline{surface} (\inserted{carte} cellulaire) \uline{à bord} (compacte), chose qui se présente quand on a un graphe $\Gamma=(S,\vec A,\sigma,\unsure{o|\alpha})$ pouvant avoir des \deleted{bouts libres} \uline{arêtes libres} (éléments \inserted{de $\vec A$} laissés fixes par $\sigma$), qui \deleted{\unsure{du coup}}\inserted{\unsure{pourront ainsi}} s'interpréter comme des ``arêtes orientées sans extrémité'' (n'allant nulle part). Quand on forme les disques polygonaux $\mathrm{Dis}_s$, et qu'on les recolle par des isom d'identification pour les $\vec a$ \uline{ordinaires} ($\sigma\vec a\ne\vec a$), on trouve tjrs une surface à bord.

Ici il s'agit de boucher les trous qui restent, \deleted{\unsure{fois}} une fois bouchés les ``trous quadrangulaires'' qu'on vient d'étudier (définis par des circuits d'arcs d'inertie) par les disques correspondants aux positions relatives \uline{supercritiques}, qu'il s'agit de recoller \deleted{aux} suivant les arcs de spécialisation exceptionnels correspondants aux arêtes libres. On est donc amené à étudier les \uline{circuits} formés d'\deleted{arêtes}\inserted{arcs} de spécialisation exceptionnels, et voir comment les boucher.

\deleted{\unsure{Alors}} Il \unsure{sera} plus simple peut être de s'inspirer de ce qui a déjà été fait, pour voir comment l'étendre au cas des spécialisations et pos.rel exceptionnelles ($=$ supercritiques). Tout d'abord, pour une \deleted{$\unsure{D_i}$} $L=D_i\in\Sigma$, on définit encore $\Delta_L$, en considérant $D_i$ comme p.r. p.r. à $\underline\Sigma=(X,\Sigma\setminus\{D_i\})$, ---
